\documentclass[10pt,a4paper]{amsart}
\usepackage[margin=19mm]{geometry}
\usepackage{amsmath,amssymb,mathtools}
\usepackage[scr=rsfs,cal=euler]{mathalfa}
\usepackage{xfrac}
\usepackage[unicode,hidelinks,bookmarks=false]{hyperref}
\newcommand{\ie}{i.e.\ }
\newcommand{\cf}{cf.\ }
\newcommand{\resp}{respectively\ }
\newcommand{\Subsection}{\subsection}
\providecommand{\nobreakdash}{\nobreak\hbox{-}}
\setlength{\emergencystretch}{2em}
\multlinegap0pt
\usepackage{dsfont}
\usepackage{amssymb}
\usepackage{enumerate}
\usepackage[shortlabels]{enumitem}
\newtheorem{claim}{Claim}
\title{A generalization of Erd\H{o}s-Hajnal problem on paths with equal-degree endpoints}
\author{Xiamiao Zhao, Yichen Wang, Mei Lu}
\date{Source: arXiv:2605.03825v1 (CC BY 4.0)}
\begin{document}
\maketitle

\noindent\emph{Source and licence.} A generalization of Erd\H{o}s-Hajnal problem on paths with equal-degree endpoints. Version arXiv:2605.03825v1 (CC BY 4.0), distributed by arXiv under the Creative Commons Attribution 4.0 International licence, http://creativecommons.org/licenses/by/4.0/, as recorded in the arXiv licence metadata for this exact version and shown on the abstract page https://arxiv.org/abs/2605.03825v1. Full licence notice: \texttt{licenses/arxiv-2605.03825-CC-BY-4.0.txt}.\par\smallskip

\noindent\textit{Editorial scope.} Counted unit: the article's claim y (source0.tex lines 384--386). The article's complete original proof of it is delivered with it (source0.tex lines 388--488, one \texttt{proof} environment of 101 lines). No mathematics is written, repaired or re-derived: the statement and the proof are the article's own lines, in the article's own notation, and no retained line is substituted or re-flowed.  Retained uncounted as the source-local context the proof needs: lines 196--198; lines 235--241; lines 245--247; lines 341--347; lines 349--357; lines 360--366; lines 369--373.  One declared adaptation: exactly 1 ancillary strings inside the retained spans are omitted (image-inclusion commands and, where the source repeats a label, the repeated label definitions) and no other line differs from the source; the figure environments, with their placement, captions and labels, are kept, and the package records the omitted strings in fidelity receipts bound to the affected bodies.  No retained line contains a citation, so the wrapper carries no bibliography and no bibliography entry.  The retained lines contain the article's own diagram code, which the wrapper typesets with the standard tikz packages; no image file is delivered and the package declares no asset dependency.  Editorial formatting only, in addition: an AMS \texttt{amsart} preamble that restates the article's own declarations and macros used by the retained lines, namely the environments \texttt{claim} on separate counters, together with the standard packages those lines need.  The retained environments print in this document as Claim 1 in order of appearance, whereas the article prints the same items with the numbering of its own full document.  The complete native source of the article as distributed in the arXiv e-print for this exact version is bundled unchanged as \texttt{sources/199818/source0.tex}.
\begin{equation}\label{eq: missing edges}
	|\mathcal{M}| = \binom{2n+1}{2} - e(G)\le n^2.
\end{equation}

\begin{equation}\label{eq: using c}
\begin{aligned}
	& b = 2(\beta - \alpha) - b_{12} = 2n + 2c - b_{12} = 2n-1-d, \\
	& d = 2n-1 - b = b_{12}-2c-1, \\
	& b_1 = b_2 = \beta - \alpha - b_{12} = c+n-b_{12} = n-c-1-d.
\end{aligned}
\end{equation}

\begin{claim}\label{claim: lower bound of D}
	$d > n^{0.9}$.
\end{claim}

\begin{equation}\label{eq: M[B]}
	\begin{aligned}		
		|\mathcal{M}[B \cup \{u,v\}]| & \ge 2(\beta - \alpha - b_{12})b_{12} + \binom{b_{12}}{2} + 2(\beta - \alpha - b_{12}) + (\beta-\alpha - b_{12})^2 - 8\ell b\\
		& \ge (\beta +1 - \alpha)^2 - \frac{1}{2}(b_{12}^2 + 5b_{12}) - 9\ell b\\
		& \ge (c+n)^2 - \binom{b_{12}+1}{2} - 24\ell n.
	\end{aligned}
	\end{equation}

	\begin{equation}\label{eq: M[B12,D]}
	\begin{aligned}
		|\mathcal{M}[B_{12}, D]|  &\ge b_{12}d - (d + (d - 1) + \ldots +(d-(d-\gamma+1))+ \gamma  + \ldots + \gamma ) \\
		& = b_{12}d - \frac{(d+\gamma+1)(d-\gamma)}{2} - \gamma(b_{12}-(d-\gamma)) \\
		& = \binom{b_{12}+1}{2} - \binom{\gamma}{2} -2c^2 - c - (2c+1)\gamma - b_{12}\\
		& \ge \binom{b_{12}+1}{2} - \binom{\gamma}{2} +2c^2+3c+1-2(c+1)b_{12}\\
		& \ge \binom{b_{12}+1}{2} - \binom{\gamma}{2} -2nc - 2n.
 	\end{aligned}
	\end{equation}

	\begin{equation}\label{eq: M[Y]}
	\begin{aligned}
	\left|\mathcal{M}[D] \right|\geq \left|\mathcal{M}[Y]\right| & \ge \binom{y_{12}}{2} + 2y_{12}(\gamma - y_{12}) + (\gamma-y_{12})^2 - 8\ell n \\
		& \ge \gamma^2 - \frac{1}{2}y_{12}^2 - \frac{1}{2}y_{12} - 8\ell n\\
		& \ge \binom{\gamma}{2} - 8\ell n,
	\end{aligned}
	\end{equation}

	\begin{equation}\label{eq: original extra linear term}
	\begin{aligned}
		n^2 \ge |\mathcal{M}| & \ge c^2 + n^2 - 29\ell n.
	\end{aligned}
	\end{equation}

\begin{claim}\label{claim: lower bound of gamma y}
	For some fixed large constants $C = C(\ell)$ depending on $\ell$ only, we have that $c \le C n^{0.5}$,  $y_{12} \ge \gamma - C n^{0.5}$, and $\gamma \ge d - 2C n^{0.5}$.
\end{claim}

\begin{proof}
	By (\ref{eq: original extra linear term}), we easily have that $c \le C n^{0.5}$ where $C$ is large enough.

	If $y_{12} < \gamma - C n^{0.5}$, then clearly $\gamma \ge Cn^{0.5}$, and we can get a refined bound for $|\mathcal{M}[Y]|$ from (\ref{eq: M[Y]}). That is
	\begin{equation}\label{eq: M[Y] 2}
	\begin{aligned}
		|\mathcal{M}[Y]| \ge & \gamma^2 - \frac{(\gamma - C n^{0.5})^2}{2} - \frac{\gamma - C n^{0.5}}{2} - 8\ell n \\
		& \ge \binom{\gamma}{2} + \frac{C^2}{2} n - 8\ell n.
	\end{aligned}
	\end{equation}
		Then combining (\ref{eq: M[B]}), (\ref{eq: M[B12,D]}) and (\ref{eq: M[Y] 2}), we have that
	\begin{equation}\label{eq: lb M}
	\begin{aligned}
		 |\mathcal{M}| & \ge c^2 + n^2 - 33\ell n  +\frac{C^2}{2}n.
	\end{aligned}
	\end{equation}
	Since $C$ is a large constant, we have a contradiction with (\ref{eq: missing edges}).

	Now we focus on the third statement.
	Suppose otherwise, $\gamma < d - 2C n^{0.5}$.
	Let $Z = D\setminus Y$, and define
	\begin{equation*}
	\begin{aligned}
		Z_1 &= \left\{w \in Z \mid d_Y(w) \le \frac{1}{2}y +2\ell \right\},\\
		Z_2 &= Z\setminus Z_1.
	\end{aligned}
	\end{equation*}
	By assumption, we have that $$z = d - y=d-(2\gamma-y_{12})=d-\gamma-(\gamma-y_{12}) \ge 2C n^{0.5}-C n^{0.5}=C n^{0.5}.$$
	We first claim that $z_2 \le \frac{1}{2}C n^{0.5}$.
	Suppose otherwise, $z_2 > \frac{1}{2}C n^{0.5}$.
	If $\mathcal{M}[Z_2] \ge \frac{1}{8}C n$, then we can add $\frac{1}{8}C n$ more missing edges to $\mathcal{M}[D]$ in (\ref{eq: M[Y]}). Combining (\ref{eq: M[B]}), (\ref{eq: M[B12,D]}), (\ref{eq: M[Y]}) and (\ref{eq: original extra linear term}), we have $|\mathcal{M}|> n^2$, a contradiction with (\ref{eq: missing edges}).
	Therefore, there exists an edge in $G[Z_2]$, say $w_1w_2$.
	Choose arbitrary vertices $w_3, w_4,\ldots,w_{\ell-1} \in Z_2$.
	By the definition of $Z_2$, for every $i=2,3,\ldots,\ell-2$, $w_i$ and $w_{i+1}$ have a distinct common neighbor in $Y$, say $t_i$~(see Figure~\ref{fig: Z2}).

	% Note that $y \ge \gamma \ge d - C n^{0.5} \ge n^{0.8}$ by Claim~\ref{claim: lower bound of D}.
	% Then we have that
	% \begin{equation*}
	% \begin{aligned}
	% 	y-y_{12} & = 2\gamma - 2y_{12} \le 2Cn^{0.5} \le 0.1y.
	% \end{aligned}
	% \end{equation*}
	Since $d_Y(w_1) \ge \frac{1}{2}y + 2\ell$, $w_1$ has a neighbor $t_1$ in $Y_1 \cup Y_{12}$ and similarly, $w_{\ell-1}$ has a neighbor $t_{\ell-1}$ in $Y_2 \cup Y_{12}$ such that $t_1,t_2,\ldots,t_{\ell-1}$ are distinct vertices.
	Then we have a path $u_1v_1t_1w_1w_2t_2w_3\ldots w_{\ell-1}t_{\ell-1}v_2u_2$ of length $2\ell -1$ connecting $u_1$ and $u_2$~(see Figure~\ref{fig: Z2}), a contradiction.

	\begin{figure}
		\centering
		
		\caption{An illustration of the path $u_1v_1t_1w_1w_2t_2w_3\ldots w_{\ell-1}t_{\ell-1}v_2u_2$.}\label{fig: Z2}
	\end{figure}

	Now we have that $z_2 \le\frac{1}{2} C n^{0.5}$. Then $z_1 = z - z_2 \ge \frac{1}{2} C n^{0.5}$ by the lower bound of $z$.
	By the definition of $Z_1$, we have that $|\mathcal{M}[Z_1, Y]| \ge \left\lfloor \frac{1}{2} y - 2\ell \right\rfloor z_1$.
	If $\gamma > n^{0.5}$, then $|\mathcal{M}[Z_1, Y]|  \ge \gamma z_1/4 \ge C n/8$, and we can also add this term to $\mathcal{M}[D]$, and get a contradiction  with (\ref{eq: missing edges}) by (\ref{eq: M[B]}), (\ref{eq: M[B12,D]}), (\ref{eq: M[Y]}) and (\ref{eq: original extra linear term}) when $C$ is large enough.
	Then we conclude that $\gamma \le n^{0.5}$. We complete the proof by considering two cases.

	\textbf{Case 1:} $\beta \ge 2b_{12}$.
	In this case, we partition $Z$ in another way.
	Define
	\begin{equation*}
	\begin{aligned}
		Z_1' &= \left\{w \in Z \mid d_{B}(w) \ge \frac{1}{2}b + 2\ell \right\},\\
		Z_2' &= Z\setminus Z_1'.
	\end{aligned}
	\end{equation*}

	We claim that $z_1' \le Cn^{0.5}$.
	Suppose $z_1' > Cn^{0.5}$. Then we claim that $G[Z_1']$ cannot contain any edge.
	The statement can be proved by a similar argument as before.
	Then $|\mathcal{M}[Z_1']| \ge \binom{z_1'}{2} \ge \frac{C^2}{4}n$, and we can also add this term to $\mathcal{M}[D]$, and get a contradiction   with (\ref{eq: missing edges}) by (\ref{eq: M[B]}), (\ref{eq: M[B12,D]}), (\ref{eq: M[Y]}) and (\ref{eq: original extra linear term}) when $C$ is sufficiently large.
	
	Then  $z_2' = d - \gamma - z_1' \ge n^{0.8}$ by Claim \ref{claim: lower bound of D} and $\gamma \le n^{0.5}$.
	For each $w \in Z_2'$, it has at least $\frac{1}{2}b - 2\ell$ missing edges in $B$.
	Since $b_{12} \le \beta/2$, we have that $\frac{1}{2} b - b_{12} \ge \frac{1}{4}\beta \ge \frac{1}{4}n$.
	Then $|\mathcal{M}[B\setminus B_{12}, Z_2']| \ge \frac{1}{4}n z_2' \ge \frac{1}{4}n^{1.8}$.
	We can also add this term to $\mathcal{M}[D]$, and get a contradiction   with (\ref{eq: missing edges}) by (\ref{eq: M[B]}), (\ref{eq: M[B12,D]}), (\ref{eq: M[Y]}) and (\ref{eq: original extra linear term}).


	\textbf{Case 2:} $n < \beta < 2b_{12}$.
	Let us partition $Z$ in another way.
	Define
	\begin{equation*}
	\begin{aligned}
		Z_1'' &= \{w \in Z \mid d_{B_{12}}(w) \ge 2\ell \},\\
		Z_2'' &= Z\setminus Z_1''.
	\end{aligned}
	\end{equation*}
	It is clear that $G[Z_1'']$ is $P_{2\ell-2}$-free. If $z_1'' > Cn^{0.5}$, then $|\mathcal{M}[Z_1'']|\ge 1/2C^2n-O(n^{0.5})$. We can  add this term to $\mathcal{M}[D]$, and get a contradiction   with (\ref{eq: missing edges}) by (\ref{eq: M[B]}), (\ref{eq: M[B12,D]}), (\ref{eq: M[Y]}) and (\ref{eq: original extra linear term}).
So we have that $z_1'' \le Cn^{0.5}$.
	Recall that $d = b_{12} - 2c - 1$ from (\ref{eq: using c}) and $c\le Cn^{0.5}$ by Claim \ref{claim: lower bound of gamma y}. Then we have a better estimation of $\mathcal{M}[B_{12}, D]$:
	\begin{equation}
	\begin{aligned}
		|\mathcal{M}[B_{12}, D]| & \ge z_2''(b_{12} - 2\ell)\\
		& \ge \left(d - 2\gamma - z_1'' \right) (b_{12} - 2\ell) \\
		& \ge \left(b_{12} - ( 2 + C) n^{0.5} -2c-1 \right) (b_{12} - 2\ell) \\
		& \ge \binom{b_{12}+1}{2} + b_{12}^2/2  - O(n^{1.5}) \\
		& \ge \binom{b_{12}+1}{2} + n^2/8 - O(n^{1.5}).
 	\end{aligned}
	\end{equation}
	We have an extra $n^2$-term compared to (\ref{eq: M[B12,D]}), a contradiction   with (\ref{eq: missing edges}) by combining (\ref{eq: M[B]}) and (\ref{eq: M[Y]}).
\end{proof}

\end{document}
